\begin{houseviewbox}[模型结论：H1 已锁定利润区间，H2 决定估值等级]
公司 H1 归母预告为 11--13 亿元，意味着 2026 年全年结果不再取决于“有没有修复”，而取决于 H2 能否在不牺牲现金的条件下延续修复。House 基准取 H1 12.0 亿元、H2 8.1 亿元，形成全年 20.1 亿元归母和 1.74 元 EPS；这不是将 Q2 年化，也不是把 H1 上沿当作全年常态。
\end{houseviewbox}

\section{2025A 起点：收入、毛利、净利与现金不能只选一个}

2025 年合并收入 85.43 亿元、归母净利 6.32 亿元，但经营现金流为 -5.70 亿元。锂产品收入 48.24 亿元、毛利率 12.22\%，民爆产品与爆破服务合计毛利 12.42 亿元，构成更稳定但并非无风险的利润底座。模型因此把锂盐的利润弹性与民爆/矿服的防御性分开建模，并将现金流单列为估值折价变量。

\begin{exhibitbox}[表 11-1：2025A 分部起点与 2026E 基准桥]
\begin{tabularx}{\exhibitboxwidth}{L{2.8cm}R{1.8cm}R{1.7cm}R{1.8cm}R{1.9cm}X}
\toprule
业务 & 2025A 收入 & 2025A 毛利率 & 2026E 收入 & 2026E 毛利率 & 基准假设 \\
\midrule
锂产品 & 48.24 & 12.22\% & 94.50 & 32\% & 9 万吨、10.5 万元/吨收入代理 \\
民爆产品 & 21.41 & 47.22\% & 22.00 & 44\% & 不以新增项目抬升利润 \\
爆破服务 & 11.75 & 19.67\% & 12.50 & 19\% & 温和增长且回款待验证 \\
运输 & 4.03 & 13.29\%* & 4.30 & 13.29\%* & *为派生毛利率代理 \\
合并 & 85.43 & 22.06\% & 133.30 & 32.16\% & 归母 20.1 亿元、EPS 1.74 元 \\
\bottomrule
\end{tabularx}
\end{exhibitbox}

\section{H1 到 H2：不把预告变成永续盈利}

H1 的预告区间是模型唯一的强前瞻锚。熊、基、牛分别取 11.0、12.0、13.0 亿元 H1 归母，H2 分别为 4.0、8.1、11.7 亿元。基准 H2 高于熊市并不意味着一路顺风，而是要求下半年锂业务收入约 36.1 亿元、毛利约 10.8 亿元，同时非锂业务提供约 6.3 亿元毛利。若正式 H1 的收入、毛利和现金无法支撑这个结构，应该改变模型变量，而不是保留结论并补充叙事。

\begin{exhibitbox}[表 11-2：H1 预告锚与 H2 独立预测]
\begin{tabularx}{\exhibitboxwidth}{L{3.0cm}R{2.0cm}R{2.0cm}R{2.0cm}X}
\toprule
项目（亿元，除销量） & 熊 & 基准 & 牛 & 模型含义 \\
\midrule
H1 归母净利 & 11.0 & 12.0 & 13.0 & 业绩预告下/中/上端 \\
H2 锂盐销量 & 3.0 万吨 & 3.8 万吨 & 4.5 万吨 & 不是 Q2 年化 \\
H2 锂业务收入 / 毛利 & 24.0 / 4.8 & 36.1 / 10.8 & 47.3 / 16.5 & House 量价成本桥 \\
H2 非锂毛利 & 5.8 & 6.3 & 6.9 & 民爆、服务、运输保守组合 \\
H2 归母净利 & 4.0 & 8.1 & 11.7 & 扣除期间费用、税和少数股东 \\
2026E 归母净利 & 15.0 & 20.1 & 24.7 & 可复算全年输出 \\
\bottomrule
\end{tabularx}
\end{exhibitbox}

\section{2027 归一化：把周期回落写进主估值方法}

对于强周期业务，估值不应只用盈利高点。2027E 基准把锂销量设为 8.8 万吨、实现收入代理 8.5 万元/吨、毛利率 24\%，归母 14.26 亿元、EPS 1.24 元；即低于 2026 基准的 20.1 亿元与 1.74 元。主 PE 法对 2026E 和 2027E 分别赋 70\% 与 30\% 权重，目的不是预测精确拐点，而是限制把 H1 利润预告永久化的倾向。

\begin{exhibitbox}[表 11-3：2027E 周期归一化，不以 2026 高点外推]
\begin{tabularx}{\exhibitboxwidth}{L{3.2cm}R{2.0cm}R{2.0cm}R{2.0cm}X}
\toprule
项目 & 熊 & 基准 & 牛 & 纪律 \\
\midrule
锂盐销量 & 8.0 万吨 & 8.8 万吨 & 9.5 万吨 & 均低于设计口径 \\
实现收入代理 & 7.8 万元/吨 & 8.5 万元/吨 & 9.4 万元/吨 & 周期回归的 House 假设 \\
锂业务毛利率 & 20\% & 24\% & 28\% & 不加入未验证成本优势 \\
合并收入 & 100.2 亿元 & 114.2 亿元 & 129.9 亿元 & 非锂业务仅温和增长 \\
归母 / EPS & 11.70 / 1.02 & 14.26 / 1.24 & 18.60 / 1.61 & 估值中的 30\% 权重输入 \\
\bottomrule
\end{tabularx}
\end{exhibitbox}

\section{利润质量压力测试：现金才是模型的第二条腿}

2026Q1 应收账款较年初增加 6.54 亿元、存货增加 1.77 亿元，经营现金流为 -4.31 亿元。基准模型仍只给全年 2.0 亿元 OCF、3.5 亿元资本开支和 -1.5 亿元自由现金流；牛情景的 6.0 亿元 OCF必须有回款和营运资本改善支撑。即使 EPS 达标、现金未达标，基准 PE 也不应上调。

\sourcenote{公司 2025 年报、2026Q1 报告与 2026H1 业绩预告；House 分部预测桥、增长模型和机器可读输入。}
